\documentclass[10pt,a4paper]{amsart}
\usepackage[margin=19mm]{geometry}
\usepackage{amsmath,amssymb,mathtools}
\usepackage[scr=rsfs,cal=euler]{mathalfa}
\usepackage{xfrac}
\usepackage[unicode,hidelinks,bookmarks=false]{hyperref}
\newcommand{\ie}{i.e.\ }
\newcommand{\cf}{cf.\ }
\newcommand{\resp}{respectively\ }
\newcommand{\Subsection}{\subsection}
\providecommand{\nobreakdash}{\nobreak\hbox{-}}
\setlength{\emergencystretch}{2em}
\multlinegap0pt
\usepackage{amssymb}
\usepackage[all]{xy}
\usepackage{mathtools}
\usepackage{xcolor}
\usepackage[normalem]{ulem}
\usepackage{nicefrac}
\usepackage{xfrac}
\usepackage[shortlabels]{enumitem}
\newtheorem{lemma}{Lemma}
\newcommand\bA{{\mathbb A}}
\newcommand\bC{{\mathbb C}}
\newcommand\bK{{\mathbb K}}
\newcommand\bR{{\mathbb R}}
\newcommand\bZ{{\mathbb Z}}
\newcommand\cF{{\mathcal{F}}}
\newcommand\cT{{\mathcal{T}}}
\newcommand{\diff}{\textnormal{d}}
\newcommand{\mb}{\mathbb}
\newcommand{\on}{\operatorname}
\title{On toric and toroidal foliations}
\author{Chih-Wei Chang, Yen-An Chen}
\date{Source: arXiv:2308.05053v4 (CC BY 4.0)}
\begin{document}
\maketitle

\noindent\emph{Source and licence.} On toric and toroidal foliations. Version arXiv:2308.05053v4 (CC BY 4.0), distributed by arXiv under the Creative Commons Attribution 4.0 International licence, http://creativecommons.org/licenses/by/4.0/, as recorded in the arXiv licence metadata for this exact version and shown on the abstract page https://arxiv.org/abs/2308.05053v4. Full licence notice: \texttt{licenses/arxiv-2308.05053-CC-BY-4.0.txt}.\par\smallskip

\noindent\textit{Editorial scope.} Counted unit: the article's lemma lem:exist\_semi\_local\_model (source0.tex lines 994--1001). The article's complete original proof of it is delivered with it (source0.tex lines 1003--1039, one \texttt{proof} environment of 37 lines). No mathematics is written, repaired or re-derived: the statement and the proof are the article's own lines, in the article's own notation, and no retained line is substituted or re-flowed.  Retained uncounted as the source-local context the proof needs: lines 926--932; lines 954--957; lines 963--966.  Nothing was omitted from any retained span, so the package carries no omission record.  No retained line contains a citation, so the wrapper carries no bibliography and no bibliography entry.  No retained line contains a figure environment, an image-inclusion command or drawing code, so no image is delivered and the package declares no asset dependency.  Editorial formatting only, in addition: an AMS \texttt{amsart} preamble that restates the article's own declarations and macros used by the retained lines, namely the environments \texttt{lemma} on separate counters, together with the standard packages those lines need.  The retained environments print in this document as Lemma 1 in order of appearance, whereas the article prints the same items with the numbering of its own full document.  The complete native source of the article as distributed in the arXiv e-print for this exact version is bundled unchanged as \texttt{sources/145426/source0.tex}.
\begin{lemma}\label{lem_property_dagger}
\begin{enumerate}
    \item  Let $N$ be a lattice, $\sigma$ be a rational, strongly convex, polyhedral cone in $N\otimes\bR$, and $W\subseteq N\otimes\bK$ be a proper $\bK$-vector subspace where $\bK=\bR$ or $\bC$. Then $(\Sigma_\sigma,W)$ satisfies the condition $(\dagger)$ if and only if $\mb R(W\cap N)\cap \sigma$ is a face of $\sigma$.
    \item Let $(\Delta,W)$ be an extended complex satisfying the condition $(\dagger)$. 
    If $\Delta'$ is a subdivision of $\Delta$, then the induced extended complex $(\Delta',W')$ also satisfies the condition $(\dagger)$. 
\end{enumerate}
\end{lemma}

\begin{lemma}\label{lem:local_isomorphism_of_affine_toric_variety}
    Let $N$ be a lattice of rank $n$, and let $U_{\sigma,N}$ be an affine toric variety. 
    If $\{m_1,\ldots,m_n\}$ is a basis of $M:=\on{Hom}(N,\bZ)$, then for any $c_1,\ldots,c_n\in\bC^*$, $x_i:=\chi^{m_i}\mapsto c_i x_i$ defines an automorphism on $T_N$ that extends to an automorphism on $U_{\sigma,N}$. 
\end{lemma}

\begin{lemma}\label{lem:minimal_intersection}
    Let $\mathbf 0\in Z\subseteq U\subseteq\bA^n=\{(x_1,\ldots,x_n)\}$ be a analytic subvariety, where $U$ is an analytic open subset, and $Z$ is smooth at $\mathbf 0$. Let $\cF$ be the foliation on $U$ defined by $\diff x_{r+1}\wedge\cdots\wedge\diff x_n$, and suppose
    $\dim_\bC \cT_{Z}(z)\cap\cF(z), z\in Z$, locally attains its minimum at $\mathbf 0$. Then there exist analytic functions $y_1,\ldots,y_n$ defined locally near $\mathbf 0$, such that $\bC[\![x_1,\ldots,x_n]\!]=\bC[\![y_1,\ldots,y_n]\!]$, $\cF$ is defined by $\diff y_{r+1}\wedge\cdots\wedge\diff y_n$, and $Z=\{y_i=0\mid i\in I\}$ in an analytic neighborhood of $\mathbf 0$ for some $I\subseteq\{1,\ldots,n\}$.
\end{lemma}

\begin{lemma}\label{lem:exist_semi_local_model}
    Let $\cF$ be a strict toroidal foliation on $X$ with extended complex $(\Delta_X,W)$. Suppose $x\in O(\tau)$ for some $\tau\in\Delta_X$.
    \begin{enumerate}
        \item There exists a semi-local model $(U_{\sigma,N},W_N)$ of $\cF$ at $x$ such that $\bR(W_{N}\cap N\cap \sigma)=W(\tau)+\bR\tilde\tau$, where $\tilde\tau\preceq\sigma$ is the face generated by all the rays $\rho\preceq\sigma$ contained in $W_N$ and $\rho\nsubseteq N_\tau\otimes\bR$. In particular, if $(\Sigma_{\tau,N_\tau}, W(\tau))$ satisfies the condition $(\dagger)$, then so does $(\Sigma_{\sigma,N},W_N)$.

        \item If $Z\subseteq V(\tau)$ is a subvariety that intersects $O(\tau)$, then we can choose a point $x\in Z\cap O(\tau)$ and a semi-local model $(U_{\sigma,N},W_N)$ of $\cF$ at $x$ such that, in addition to the statement in (1), $Z$ formally locally corresponds to $V_{\tau_Z}$ for some $\tau\preceq\tau_Z\preceq \sigma$.
    \end{enumerate}
\end{lemma}

\begin{proof}
    We first prove (1). By considering a local model $(U_{\tau,\widetilde N},W_{\widetilde N})$ of $\cF$ at $x$, we may replace $X$ by $U_{\tau,\widetilde N}$. Choosing a sublattice $N_0\subseteq \widetilde N$ such that $\widetilde N=N_\tau\oplus N_0$, we get $U_{\tau,\widetilde N}\simeq U_{\tau,N_\tau}\times T_{N_0}$, and we may assume that $x=(\gamma_\tau,\mathbf 1)\in U_{\tau,N_\tau}\times T_{N_0}$. Choose bases $\{m_1,\ldots,m_k\}$ and $\{m_1',\ldots,m_\ell'\}$ of $\on{Hom} (N_\tau,\bZ)$ and $\on{Hom} (N_0,\bZ)$, respectively, and write $x_i=\chi^{m_i},y_j=\chi^{m'_j}$. By expressing $W_{\widetilde N}$ as an intersection of hyperplanes of $\widetilde N\otimes\bC$, we may assume that $\cF_{W_{\widetilde N}}$ is defined by $\omega=\omega_1\wedge\cdots\wedge\omega_c$, where $c=\dim X-\on{rank}\cF$ and $\omega_i=\sum_{j=1}^ka_{ij}\frac{\diff x_j}{x_j}+\sum_{j=1}^{\ell}b_{ij}\frac{\diff y_{j}}{y_{j}}$. Let $A$ be the matrix defined by $A_{ij}=a_{ij}$ if $1\leq j\leq k$, and $A_{ij}=b_{i,j-k}$ if $k+1\leq j\leq k+\ell$. By performing row operations, possibly re-indexing the variables, and modifying $\omega_i$ accordingly, we may assume that  
    \begin{equation}\tag{*}
        A =
        \begin{bmatrix}
        B & \textbf{0} & *\\
        \mathbf{0} & C & *
        \end{bmatrix},
    \end{equation}
    where
    \begin{itemize}
        \item  $B$ is an $ m \times k $ matrix in row-echelon form of full rank $m$ with pivots in the first $m$ columns, and
        \item  $C$ is a $(c-m) \times (c-m)$ diagonal square matrix of full rank.
    \end{itemize} 
    We allow $m$ to be $0$ or $c$, and choose all the pivots to be $1$.
    Here we use the fact $\on{rank}(A)=c$.
    
    Note that centered at $x$, $\cF_{W_{\widetilde N}}$ is defined by the wedge product of $\omega_i=\sum_{j=1}^ka_{ij}\frac{\diff x_j}{x_j}+\sum_{j=1}^{\ell}b_{ij}\frac{\diff z_{j}}{z_{j}+1}$, where $z_j=y_j-1$. Assume (*), and consider the following analytic functions:
    \begin{align*}
    &x_i'=x_i(z_1+1)^{b_{i,1}}\cdots(z_\ell+1)^{b_{i,\ell}},\ 1\leq i\leq m,\\
    &x_{i}'=x_{i},\ m+1\leq i\leq k,\\
    &z'_i=(z_1+1)^{b_{m+i,1}}\cdots (z_\ell+1)^{b_{m+i,\ell}}-1,\ 1\leq i\leq c-m-1,\\
    &z'_{i}=z_{i},\ c-m\leq i\leq \ell.
    \end{align*}
    Then $(x_1',\ldots,x_k',z_1',\ldots,z_\ell')$ defines a biholomorphic map from an analytic neighborhood of $x$ to an analytic neighborhood of $(\gamma_{\tau},0)\in U_{\tau,N_\tau}\times\bA^\ell$ by Lemma~\ref{lem:local_isomorphism_of_affine_toric_variety}. Write $\bA^\ell=U_{\sigma_1,N_1}$ through the coordinates $z'_i$, and let $N=N_\tau\oplus N_1$, $\sigma=\tau\times\sigma_1$. Then $\cF_{W_{\widetilde N}}$ maps to the toric foliation $\cF_{W_N}$ on $U_{\sigma,N}$ defined by $\omega'=\omega'_1\wedge\cdots\wedge\omega_c'$, where $\omega'_i=\sum_{j=i}^ka_{ij}\frac{\diff x_j'}{x_j'}$ for $1\leq i\leq m$, $\omega'_i=\diff z'_{i-m}$ for $m+1\leq i\leq c$. Denote by $\widetilde\tau$ the face of $\sigma$ generated by all the rays $\rho\preceq\{0\}\times\sigma_1$ contained in $W_N$.
    Then one can check that $W_N\cap N=W_N\cap N_\tau+W_N\cap N_1$, and $W_N\cap N_1=N\cap\bR\widetilde\tau$. Since $W(\tau) = \bR(W_N\cap N_\tau \cap \tau)$, the assertion $\bR(W_{N}\cap N\cap\sigma)=\bR(W_{N}\cap N\cap(\tau\times\widetilde\sigma))=W(\tau)+\bR\widetilde\tau$ follows.

    Suppose $(\Sigma_{\tau,N_\tau},W(\tau))$ satisfies the condition $(\dagger)$. 
    Then by Lemma~\ref{lem_property_dagger}(1), we have $\bR(W(\tau)\cap N_\tau) \cap \tau$ is a face of $\tau$. 
    Note that 
    \[\bR(W_N\cap N)\cap\sigma=\bR(W_N\cap N\cap\sigma)\cap\sigma=(W(\tau)+\bR\widetilde\tau)\cap(\tau\times\widetilde\sigma)=(W(\tau)\cap\tau)\times \widetilde\tau.\]
    Since $W(\tau)$ is generated by lattice points, we have $W(\tau)\cap \tau=\tau_0$ for some $\tau_0\preceq\tau$. Hence $\bR(W_N\cap N)\cap\sigma=\tau_0\times\widetilde\tau\preceq\sigma$,
    and therefore $(\Sigma_{\sigma,N},W_N)$ satisfies the condition $(\dagger)$ again by Lemma~\ref{lem_property_dagger}(1).

    Now we prove (2). 
    Write $\overline\cF=\cF\vert_{V(\tau)}$, and choose $x\in O(\tau)$ to be a point such that $\dim_\bC(\overline\cF(z)\cap \cT_Z(z))$ locally attains its minimum at $x$. Assume (*), and choose a semi-local $(U_{\sigma,N},W_N)$ of $\cF$ at $x$ as in the proof of Lemma~\ref{lem:exist_semi_local_model}(1). Then $(U_{\sigma_1,N_1},\bC\widetilde\tau)$ is a semi-local model of $\overline\cF$ at $x$. Choose $\widetilde Z \subseteq U_{\sigma_1,N_1}$ defined analytically locally near $\gamma_{\sigma_1}$, such that $\widetilde Z$ maps to $Z$ formally locally near $\gamma_{\sigma_1}$. By applying Lemma~\ref{lem:minimal_intersection}, we get the analytic functions $z''_1,\ldots,z''_\ell$ on $U_{\sigma_1,N_1}$ such that $\cF_{\bC \widetilde\tau}$ is defined by $\diff z''_1\wedge\cdots\wedge\diff z''_{c-m}$, and $\widetilde Z=\{z''_i=0\mid i\in I\}$ near $\gamma_{\sigma_1}$ for some $I\subseteq\{1,\ldots,c-m\}$. This gives a toric foliation $\cF_{\bC\widehat\tau}$ on $U_{\sigma_2,N_2}$ for some $\widehat\tau\preceq\sigma_2$. Replacing $N_1$ by $N_2$ and  $(U_{\sigma_1,N_1},\bC\widetilde\tau)$ by $(U_{\sigma_2,N_2},\bC\widehat\tau)$, we get the semi-local model as required.
\end{proof}

\end{document}
